% This is text associated with the open textbook "Open Electromagnetics"
% See immediately below for author, date
% This work is licensed under the Creative Commons Attribution-ShareAlike 4.0 International License. (CC BY-SA 4.0) 
% To view a copy of the license, visit http://creativecommons.org/licenses/by-sa/4.0/

%ORB TITLE Electrostatics
%ORB AUTHOR 0        # S.W. Ellingson, Virginia Tech (ellingson@vt.edu)
%ORB DATE 20180803   # yyyymmdd
%ORB ABSTRACT Electrostatics 

\label{m0116_Electrostatics_Preface} % <-- Must match name of this file and module directory name.  Don't mess with this!
\index{electrostatics (description)|(}

\emph{Electrostatics} is the theory of the electric field in conditions in which its behavior is independent of magnetic fields, including
%
\begin{itemize}
\item The electric field associated with fixed distributions of electric charge
\item \emph{Capacitance} (the ability of a structure to store energy in an electric field)
\item The \emph{energy} associated with the electrostatic field
\item \emph{Steady current} induced in a conducting material in the presence of an electrostatic field (essentially, Ohm's Law)
\end{itemize}
%
The term ``static'' refers to the fact that these aspects of electromagnetic theory can be developed by assuming sources are time-invariant; we might say that electrostatics is the study of the electric field at DC.  However, many aspects of electrostatics are relevant to AC, radio frequency, and higher-frequency applications as well.
 
%%%%%%%%%%%%%%%%%%%%%%%%%%%%%%%%%%%%%%%%

{\bf Additional Reading:}
\begin{itemize}
\item ``\href{https://en.wikipedia.org/wiki/Electrostatics}{Electrostatics}'' on Wikipedia.
\end{itemize}

\index{electrostatics (description)|)}


\vspace{1in}
